\begin{frame}{자가진단 답 (한 줄) — 핵심 10문항}
  \scriptsize
  \begin{enumerate}
    \item \textbf{1} — 목표와 추정치가 서로를 따라다니며 발산하기
      때문(9.2절 보스팅 불안정의 방치 버전).
    \item \textbf{3} — $\delta = r + \gamma Q(s',a') - Q(s,a)$ —
      최적을 추구하면 $a'$ 자리에 $\max$.
    \item \textbf{5} — baseline $b(s)$가 $a$에 무관하므로
      $\mathbb{E}_a[\nabla \log \pi(a|s)\, b(s)] = b(s)\nabla
      \sum_a \pi(a|s) = 0$ — \textbf{분산만 줄고 편향은 0}.
    \item \textbf{6} — $r_t(\theta)$가 1에서 크게 벗어나 한
      업데이트에 정책이 ``멀리 점프''해 좋은 정책을 \textbf{망가뜨리는
      것} — 클리핑은 점프 반경을 제한.
    \item \textbf{8} — $\lambda=0$이면 TD(0)(편향 있음, 분산 작음),
      $\lambda\to 1$이면 (유한 에피소드에서) MC 리턴(편향 없음,
      분산 큼) — 7.2의 U자 스펙트럼.
    \item \textbf{9} — 전문가 데이터 \textbf{분포 밖} 상태에 들면 그
      행동이 신뢰 불가 $\to$ 다음 상태가 더 벗어난 분포로 $\to$
      오차가 \textbf{스텝마다 곱해짐}(covariate shift).
    \item \textbf{11} — 사람의 \textbf{선호}(두 응답 중 어느 것이
      나은지) — 선호 데이터로 \textbf{보상 모델}을 먼저 학습하고, 그
      보상으로 PPO.
    \item \textbf{13} — \textbf{systematic}한 차이이므로 시뮬레이션
      데이터를 늘려도 \textbf{사라지지 않}음 — ML1 Ch04
      bias-variance에서 ``편이는 데이터로 안 줄고''의 RL 버전.
    \item \textbf{15} — MuJoCo: 접촉·마찰 \textbf{정밀도}, Isaac: GPU
      \textbf{병렬화}(수천 환경 동시) — 정밀도 대 규모.
    \item \textbf{20} — ML1은 \emph{다른이(데이터 수집자)가 만든
      데이터와 정답}을 줘 맞히게 하고, ML2는 \emph{에이전트가
      행동해서 데이터를 스스로 만들고 정답 없이 보상만으로} 배운다 —
      ``\textbf{무엇을 데이터로 갖고, 무엇을 최적화하는가}''의
      답이 바뀌는 것.
  \end{enumerate}
\end{frame}
